\documentclass[a4paper,12pt]{article}

% -------------------------------------------------
% Кодировка и язык
% -------------------------------------------------
\usepackage[T2A]{fontenc}
\usepackage[utf8]{inputenc}
\usepackage[russian]{babel}

% -------------------------------------------------
% Математика
% -------------------------------------------------
\usepackage{amsmath,amssymb,mathtools}
\usepackage{icomma}

% -------------------------------------------------
% Типографика
% -------------------------------------------------
\usepackage{helvet}
\renewcommand{\familydefault}{\sfdefault}
\usepackage{microtype}
\usepackage{setspace}
\usepackage{parskip}
\usepackage{enumitem}

% -------------------------------------------------
% Страница
% -------------------------------------------------
\usepackage[
  left=22mm,
  right=22mm,
  top=22mm,
  bottom=24mm
]{geometry}

% -------------------------------------------------
% Графика
% -------------------------------------------------
\usepackage{xcolor}
\definecolor{accent}{HTML}{008080}
\definecolor{darkaccent}{HTML}{004D4D}
\definecolor{textgray}{HTML}{333333}
\definecolor{softgray}{HTML}{F3F5F5}

\usepackage{tikz}
\usetikzlibrary{calc}

\usepackage{tkz-euclide}

\usepackage{pgfplots}
\pgfplotsset{compat=1.18}

% -------------------------------------------------
% Таблицы
% -------------------------------------------------
\usepackage{tabularx,booktabs}

% -------------------------------------------------
% Служебные пакеты
% -------------------------------------------------
\usepackage{hyperref}
\hypersetup{
  colorlinks=true,
  linkcolor=darkaccent,
  urlcolor=darkaccent
}

\usepackage{fancyhdr}

% -------------------------------------------------
% Общие настройки
% -------------------------------------------------
\onehalfspacing
\setlength{\parindent}{0pt}
\setlength{\headheight}{15pt}
\setlist[itemize]{leftmargin=7mm,itemsep=3pt,topsep=4pt}
\setlist[enumerate]{leftmargin=9mm,itemsep=5pt,topsep=5pt}

\pagestyle{fancy}
\fancyhf{}
\fancyhead[L]{\small\color{textgray}Чек-лист по математике}
\fancyhead[R]{\small\color{textgray}Подготовка к ОГЭ}
\fancyfoot[C]{\small\color{textgray}\thepage}
\renewcommand{\headrulewidth}{0.3pt}
\renewcommand{\footrulewidth}{0pt}

\newcommand{\topic}[1]{%
  \vspace{0.4em}
  {\Large\bfseries\color{darkaccent}#1\par}
  \vspace{0.25em}
  {\color{accent}\rule{\linewidth}{0.7pt}}
  \vspace{0.45em}
}

\newcommand{\subtopic}[1]{%
  \vspace{0.5em}
  {\large\bfseries\color{darkaccent}#1\par}
  \vspace{0.2em}
}

\newcommand{\important}[1]{%
  \par\vspace{0.3em}
  {\color{darkaccent}\bfseries Важно.} #1
  \par\vspace{0.3em}
}

\newcommand{\checkliststart}{%
  \begin{itemize}[label={\(\square\)},leftmargin=8mm]
}
\newcommand{\checklistend}{\end{itemize}}

\begin{document}

% =================================================
% ТИТУЛЬНЫЙ ЛИСТ
% =================================================
\begin{titlepage}
\thispagestyle{empty}

\begin{center}

\vspace*{30mm}

{\color{darkaccent}\rule{0.75\linewidth}{1pt}}

\vspace{10mm}

{\Huge\bfseries Чек-лист}

\vspace{6mm}

{\LARGE\bfseries
Входная диагностика и линейная функция}

\vspace{4mm}

{\Large
Ключевые опоры для подготовки к ОГЭ}

\vspace{10mm}

{\color{accent}\rule{0.48\linewidth}{0.8pt}}

\vfill

{\large
Лёвин Артём Александрович\\[2mm]
эксклюзивно для Марины
}

\vspace{10mm}

{\large \today}

\vspace{24mm}

\end{center}

\begin{tikzpicture}[remember picture,overlay]
  \draw[
    accent!55,
    line width=1.2pt
  ]
  ([xshift=18mm,yshift=17mm]current page.south west)
  .. controls
  ([xshift=55mm,yshift=31mm]current page.south west)
  and
  ([xshift=-58mm,yshift=8mm]current page.south east)
  ..
  ([xshift=-18mm,yshift=20mm]current page.south east);
\end{tikzpicture}

\end{titlepage}

% =================================================
% ОСНОВНОЙ ЧЕК-ЛИСТ
% =================================================

\topic{1. Алгебраические основы: что обязательно помнить}

\subtopic{1.1. Квадратный корень}

Для положительных чисел полезно искать под корнем полный квадрат:
\[
\sqrt{ab}=\sqrt a\,\sqrt b,
\qquad a\ge0,\quad b\ge0.
\]

Например,
\[
\sqrt{27}
=
\sqrt{9\cdot3}
=
3\sqrt3.
\]

\checkliststart
\item Я умею находить под корнем множитель, являющийся полным квадратом.
\item Я помню, что \(\sqrt{a^2}=|a|\).
\item Если под корнем находится выражение с переменной, проверяю условие
\[
g(x)\ge0
\]
для корня \(\sqrt{g(x)}\).
\checklistend

\subtopic{1.2. Отрицательная степень}

Для \(a\ne0\):
\[
a^{-n}=\frac{1}{a^n}.
\]

Например,
\[
2^{-3}
=
\frac1{2^3}
=
\frac18
=
0{,}125.
\]

\important{Отрицательная степень меняет расположение основания в дроби. Само значение степени остаётся положительным показателем.}

\subtopic{1.3. Формулы сокращённого умножения}

На занятии особенно важны две формулы.

\[
a^2-b^2=(a-b)(a+b)
\]
--- \textbf{разность квадратов}.

Например,
\[
9x^2-4
=
(3x)^2-2^2
=
(3x-2)(3x+2).
\]

Вторая формула:
\[
(a-b)^2=a^2-2ab+b^2.
\]

Например,
\[
x^2-4x+4
=
x^2-2\cdot x\cdot2+2^2
=
(x-2)^2.
\]

\checkliststart
\item Я различаю разность квадратов и квадрат разности.
\item Перед применением формулы представляю каждую часть как квадрат.
\item После разложения могу выполнить обратную проверку раскрытием скобок.
\checklistend

% =================================================
\topic{2. ОДЗ и алгебраические дроби}

Если переменная находится в знаменателе, знаменатель должен отличаться от нуля.

Например, для выражения
\[
\frac{x}{x+5}-\frac4{x-2}
\]
получаем ограничения:
\[
x+5\ne0,
\qquad
x-2\ne0.
\]

Следовательно,
\[
\boxed{x\ne-5,\qquad x\ne2}.
\]

\subtopic{Алгоритм}

\begin{enumerate}
\item Найти все знаменатели.
\item Каждый знаменатель приравнять к нулю.
\item Полученные значения исключить.
\item Записать ОДЗ до основных преобразований.
\item Сохранить найденные ограничения до конца решения.
\end{enumerate}

\important{Сокращение дроби не отменяет исходные ограничения. Если значение было запрещено первоначальным знаменателем, оно остаётся запрещённым.}

\checkliststart
\item Я проверяю знаменатели на ноль.
\item Я записываю ОДЗ до преобразований.
\item Я сверяю найденные корни с ОДЗ перед записью ответа.
\checklistend

% =================================================
\topic{3. Неравенства и числовая прямая}

\subtopic{3.1. Деление на отрицательное число}

При умножении или делении обеих частей неравенства на отрицательное число знак неравенства меняется:

\[
-2x>4
\]

делим на \(-2\):

\[
x<-2.
\]

Это одно из ключевых правил работы с неравенствами.

\subtopic{3.2. Строгое и нестрогое неравенство}

\[
x\ge1
\]
означает: подходят \(1\) и все числа больше \(1\).

В интервальной форме:
\[
[1;+\infty).
\]

\[
x>1
\]
означает: подходят только числа больше \(1\).

В интервальной форме:
\[
(1;+\infty).
\]

Аналогично:
\[
x\le1
\quad\Longleftrightarrow\quad
(-\infty;1],
\]
\[
x<1
\quad\Longleftrightarrow\quad
(-\infty;1).
\]

\begin{center}
\begin{tikzpicture}[x=0.65cm,y=1cm]
  \draw[->,line width=0.8pt] (-9,0)--(11,0);
  \draw (-7,0.12)--(-7,-0.12);
  \draw (9,0.12)--(9,-0.12);
  \node[below] at (-7,-0.12) {$-7$};
  \node[below] at (9,-0.12) {$9$};

  \draw[accent,line width=2pt] (-7,0)--(9,0);
  \draw[accent,line width=1.2pt,fill=white] (-7,0) circle (2.4pt);
  \fill[accent] (9,0) circle (2.5pt);

  \node[above] at (1,0.35) {$(-7;9]$};
\end{tikzpicture}
\end{center}

Например,
\[
\begin{cases}
x>-7,\\
x\le9
\end{cases}
\]
имеет решение
\[
\boxed{x\in(-7;9]}.
\]

\subtopic{3.3. Система неравенств}

Решение системы --- \textbf{пересечение} всех полученных множеств.

\checkliststart
\item Я меняю знак неравенства при делении на отрицательное число.
\item Для \(<\) и \(>\) использую выколотую точку.
\item Для \(\le\) и \(\ge\) использую закрашенную точку.
\item В системе ищу общую часть всех промежутков.
\item У бесконечности всегда ставлю круглую скобку.
\checklistend

% =================================================
\topic{4. Геометрическое оформление и подобие треугольников}

\subtopic{4.1. Оформление решения}

В геометрии важно явно обозначать введённые объекты.

Если введена переменная \(x\), следует написать, какой отрезок она обозначает, например
\[
BE=x.
\]

Углы желательно указывать тремя буквами:
\[
\angle DAE,\qquad \angle ABC.
\]

Средняя буква --- вершина угла.

\important{При доказательстве равенства углов указывайте причину: общий угол, соответственные углы при параллельных прямых, вертикальные углы и т.\,д.}

\subtopic{4.2. Подобие по двум углам}

Если две пары углов двух треугольников соответственно равны, треугольники подобны.

Рассмотрим ситуацию
\[
DE\parallel BC.
\]

\begin{center}
\begin{tikzpicture}[scale=1.0]
  \tkzSetUpLine[line width=1pt,color=accent]
  \tkzSetUpPoint[color=accent,fill=accent,size=3]
  \tkzSetUpLabel[color=black,font=\small]

  \tkzDefPoint(0,4){A}
  \tkzDefPoint(-3,0){B}
  \tkzDefPoint(3,0){C}
  \tkzDefPoint(-1.8,1.6){D}
  \tkzDefPoint(1.8,1.6){E}

  \tkzDrawSegments(A,B A,C B,C D,E)
  \tkzDrawPoints(A,B,C,D,E)

  \tkzLabelPoints[above](A)
  \tkzLabelPoints[below left](B)
  \tkzLabelPoints[below right](C)
  \tkzLabelPoints[left](D)
  \tkzLabelPoints[right](E)
\end{tikzpicture}
\end{center}

Имеем:
\[
\angle DAE=\angle BAC
\]
как общий угол при вершине \(A\),

\[
\angle ADE=\angle ABC
\]
как соответственные углы при
\[
DE\parallel BC.
\]

Следовательно,
\[
\triangle ADE\sim\triangle ABC.
\]

Из подобия:
\[
\boxed{
\frac{DE}{BC}
=
\frac{AE}{AC}
=
\frac{AD}{AB}
}.
\]

Соответственные стороны лежат напротив равных углов.

\subtopic{Пример}

Пусть
\[
AD=6,\qquad DB=3,\qquad AE=8,\qquad EC=x.
\]

Тогда
\[
AB=9,\qquad AC=8+x.
\]

Используем отношение соответственных сторон:
\[
\frac{AE}{AC}
=
\frac{AD}{AB}.
\]

Получаем
\[
\frac8{8+x}=\frac69.
\]

Перемножаем крест-накрест:
\[
72=48+6x,
\]
\[
6x=24,
\]
\[
\boxed{x=4}.
\]

\checkliststart
\item Я умею доказать подобие по двум углам.
\item Я поясняю происхождение равных углов.
\item Я правильно сопоставляю соответствующие стороны.
\item Я записываю пропорцию в одном и том же порядке для обоих треугольников.
\checklistend

% =================================================
\topic{5. Центральный и вписанный углы}

Пусть \(O\) --- центр окружности.

\subtopic{5.1. Центральный угол}

Центральный угол равен градусной мере дуги, на которую он опирается:
\[
\angle TON=\overset{\frown}{TN}.
\]

\subtopic{5.2. Вписанный угол}

Вписанный угол равен половине градусной меры дуги, на которую он опирается:
\[
\angle TKN
=
\frac12\overset{\frown}{TN}.
\]

Если центральный и вписанный углы опираются на одну дугу, то
\[
\boxed{
\angle TKN=\frac12\angle TON
}.
\]

Например,
\[
\angle TON=30^\circ
\]
даёт
\[
\angle TKN=15^\circ.
\]

\subtopic{5.3. Угол, опирающийся на диаметр}

Если \(TK\) --- диаметр окружности, а \(N\) лежит на окружности, то
\[
\angle TNK=90^\circ.
\]

\begin{center}
\begin{tikzpicture}[scale=1.0]
  \tkzSetUpLine[line width=1pt,color=accent]
  \tkzSetUpPoint[color=accent,fill=accent,size=3]
  \tkzSetUpLabel[color=black,font=\small]

  \tkzDefPoint(0,0){O}
  \tkzDefPoint(-2.5,0){T}
  \tkzDefPoint(2.5,0){K}
  \tkzDefPoint(0.7,2.4){N}

  \tkzDrawCircle(O,T)
  \tkzDrawSegments(T,K T,N N,K O,N)
  \tkzDrawPoints(O,T,K,N)
  \tkzMarkRightAngle[size=0.25](T,N,K)

  \tkzLabelPoints[below](O)
  \tkzLabelPoints[left](T)
  \tkzLabelPoints[right](K)
  \tkzLabelPoints[above](N)
\end{tikzpicture}
\end{center}

Причина:
\[
\overset{\frown}{TK}=180^\circ,
\]
поэтому вписанный угол равен
\[
\angle TNK
=
\frac12\cdot180^\circ
=
90^\circ.
\]

\checkliststart
\item Я отличаю центральный угол от вписанного.
\item Я проверяю, на какую дугу опирается угол.
\item Я помню связь: вписанный угол вдвое меньше центрального при общей дуге.
\item Я узнаю прямой угол, опирающийся на диаметр.
\checklistend

% =================================================
\topic{6. Основная тема: линейная функция}

Линейная функция задаётся формулой
\[
\boxed{y=kx+b}.
\]

Её график --- прямая.

\subtopic{6.1. Основные термины}

\[
x
\]
--- независимая переменная; на координатной плоскости её значение называется \textbf{абсциссой}.

\[
y
\]
--- значение функции; на координатной плоскости его значение называется \textbf{ординатой}.

\[
k
\]
--- \textbf{угловой коэффициент}.

\[
b
\]
--- \textbf{свободный коэффициент}.

\checkliststart
\item Я знаю общий вид \(y=kx+b\).
\item Я различаю абсциссу \(x\) и ординату \(y\).
\item Я знаю смысл коэффициентов \(k\) и \(b\).
\checklistend

% =================================================
\topic{7. Что показывает коэффициент \(k\)}

Знак \(k\) определяет направление прямой.

\[
k>0
\quad\Longrightarrow\quad
\text{функция возрастает}.
\]

При движении слева направо график идёт вверх.

\[
k<0
\quad\Longrightarrow\quad
\text{функция убывает}.
\]

При движении слева направо график идёт вниз.

При
\[
k=0
\]
получаем
\[
y=b,
\]
то есть горизонтальную прямую.

\begin{figure}[ht]
\centering
\begin{tikzpicture}
\begin{axis}[
  width=0.84\linewidth,
  height=7.8cm,
  axis lines=middle,
  unit vector ratio=1 1,
  xmin=-5,
  xmax=6,
  ymin=-5,
  ymax=6,
  xtick={-4,-2,0,2,4},
  ytick={-4,-2,0,2,4},
  xlabel={$x$},
  ylabel={$y$},
  samples=220,
  clip=false,
  grid=both,
  minor tick num=1,
  grid style={gray!15},
  legend style={
    draw=none,
    fill=none,
    at={(0.98,0.98)},
    anchor=north east,
    font=\small
  }
]
\addplot[
  accent,
  very thick,
  domain=-5:6
] {0.7*x-2};
\addlegendentry{\(k>0\)}

\addplot[
  darkaccent,
  very thick,
  dashed,
  domain=-4:6
] {-0.8*x+3};
\addlegendentry{\(k<0\)}
\end{axis}
\end{tikzpicture}
\caption{Возрастающая и убывающая линейные функции}
\end{figure}

\important{Знак \(k\) можно использовать для проверки результата. Если график убывает, найденный коэффициент \(k\) должен быть отрицательным.}

% =================================================
\topic{8. Что показывает коэффициент \(b\)}

Подставим
\[
x=0.
\]

Тогда
\[
y=k\cdot0+b=b.
\]

Следовательно, график пересекает ось \(Oy\) в точке
\[
\boxed{(0;b)}.
\]

Поэтому:

\[
b>0
\]
--- пересечение с осью \(Oy\) выше начала координат;

\[
b=0
\]
--- график проходит через начало координат;

\[
b<0
\]
--- пересечение с осью \(Oy\) ниже начала координат.

\begin{table}[ht]
\centering
\caption{Четыре основных сочетания знаков \(k\) и \(b\)}
\renewcommand{\arraystretch}{1.3}
\begin{tabularx}{0.88\linewidth}{@{}c c X@{}}
\toprule
\(k\) & \(b\) & Что видно по графику \\
\midrule
\(k>0\) & \(b>0\) &
Прямая возрастает и пересекает \(Oy\) выше начала координат.\\
\(k>0\) & \(b<0\) &
Прямая возрастает и пересекает \(Oy\) ниже начала координат.\\
\(k<0\) & \(b>0\) &
Прямая убывает и пересекает \(Oy\) выше начала координат.\\
\(k<0\) & \(b<0\) &
Прямая убывает и пересекает \(Oy\) ниже начала координат.\\
\bottomrule
\end{tabularx}
\end{table}

% =================================================
\topic{9. Метод узловых точек}

\subtopic{9.1. Что такое узловая точка}

В рамках работы с графиком удобно выбирать точки, координаты которых являются целыми числами.

Например:
\[
Q(0;-2),\qquad J(3;0).
\]

Первая координата --- \(x\), вторая --- \(y\).

\subtopic{9.2. Сколько точек нужно}

Для функции
\[
y=kx+b
\]
неизвестны два коэффициента:
\[
k,\qquad b.
\]

Поэтому достаточно взять две подходящие точки графика.

\subtopic{9.3. Алгоритм метода}

\begin{enumerate}
\item Выбрать две удобные точки графика.
\item Записать их координаты в формате
\[
(x;y).
\]
\item Для каждой точки подставить \(x\) и \(y\) в
\[
y=kx+b.
\]
\item Получить систему двух уравнений.
\item Найти \(k\) и \(b\).
\item Проверить знак \(k\) по направлению графика.
\end{enumerate}

\subtopic{9.4. Пример}

Пусть прямая проходит через
\[
Q(0;-2),
\qquad
J(3;0).
\]

Подставляем первую точку:
\[
-2=k\cdot0+b,
\]
откуда
\[
b=-2.
\]

Подставляем вторую:
\[
0=3k+b.
\]

Используем \(b=-2\):
\[
0=3k-2,
\]
\[
3k=2,
\]
\[
k=\frac23.
\]

Искомая функция:
\[
\boxed{
y=\frac23x-2
}.
\]

Проверка:
\[
k=\frac23>0,
\]
поэтому прямая должна возрастать.

\subtopic{9.5. Ещё одна быстрая проверка}

Пусть прямая проходит через
\[
T(0;5),
\qquad
N(4;0).
\]

Из первой точки:
\[
b=5.
\]

Из второй:
\[
0=4k+5,
\]
поэтому
\[
k=-\frac54.
\]

Получаем:
\[
\boxed{
y=-\frac54x+5
}.
\]

Так как
\[
k<0,
\]
график убывает.

\subtopic{9.6. Формула для самопроверки}

Если известны две точки
\[
A(x_1;y_1),
\qquad
B(x_2;y_2),
\qquad
x_1\ne x_2,
\]
то
\[
\boxed{
k=\frac{y_2-y_1}{x_2-x_1}
}.
\]

После нахождения \(k\):
\[
b=y_1-kx_1.
\]

Основной рабочий метод для текущего этапа подготовки --- подстановка координат точек в \(y=kx+b\).

% =================================================
\topic{10. Финальный чек перед решением задачи}

\checkliststart
\item В корнях ищу полный квадрат и проверяю допустимость выражений с переменной.
\item В отрицательной степени использую правило \(a^{-n}=1/a^n\).
\item Узнаю разность квадратов и квадрат разности.
\item Для дробей заранее выписываю ОДЗ.
\item При делении неравенства на отрицательное число меняю знак.
\item Отличаю строгие и нестрогие границы промежутка.
\item В системе неравенств ищу пересечение.
\item В геометрии подписываю введённые точки, отрезки и переменные.
\item Равенство углов сопровождаю обоснованием.
\item Для подобных треугольников правильно сопоставляю стороны.
\item Помню связь центрального и вписанного углов.
\item Для \(y=kx+b\) определяю смысл каждого коэффициента.
\item По знаку \(k\) определяю возрастание или убывание.
\item По \(b\) определяю пересечение с осью \(Oy\).
\item Для нахождения \(k\) и \(b\) умею использовать две узловые точки.
\item После вычислений выполняю проверку здравым смыслом.
\checklistend

% =================================================
% ТРЕНИРОВОЧНЫЙ БЛОК
% =================================================
\newpage
\topic{Тренировочный блок}

\textbf{Путь А.} Задания расположены по нарастающей сложности.

\subtopic{Средний уровень}

\begin{enumerate}

\item Упростите и преобразуйте:
\[
\text{1)}\ \sqrt{75};
\qquad
\text{2)}\ 2^{-4};
\qquad
\text{3)}\ 16x^2-25;
\qquad
\text{4)}\ x^2-10x+25.
\]

\item Для выражения
\[
\frac{x+1}{x-3}-\frac5{x+2}
\]
укажите все значения \(x\), при которых выражение не имеет смысла.

\item Дана функция
\[
y=-3x+6.
\]
Определите:
\[
\text{1) знак }k;
\qquad
\text{2) возрастает или убывает функция;}
\]
\[
\text{3) точку пересечения с осью }Oy;
\qquad
\text{4) точку пересечения с осью }Ox.
\]

\end{enumerate}

\subtopic{Уровень выше среднего}

\begin{enumerate}[resume]

\item Прямая проходит через точки
\[
A(0;-4),
\qquad
B(6;2).
\]
Найдите \(k\), \(b\) и запишите уравнение прямой.

\item Прямая проходит через точки
\[
C(-2;5),
\qquad
D(4;-1).
\]
Найдите её уравнение и укажите, возрастает или убывает соответствующая функция.

\item Дана функция
\[
y=(m-2)x+5.
\]

\begin{enumerate}[label=\arabic*)]
\item При каких \(m\) функция возрастает?
\item Найдите \(m\), если график проходит через точку
\[
P(3;11).
\]
\end{enumerate}

\item Решите систему неравенств:
\[
\begin{cases}
-3x<12,\\
2x-1\le9.
\end{cases}
\]
Ответ запишите промежутком.

\item В треугольнике \(ABC\)
\[
D\in AB,\qquad E\in AC,\qquad DE\parallel BC.
\]
Известно:
\[
AD=4,\qquad DB=2,\qquad AE=6.
\]
Найдите \(EC\).

\item \(TK\) --- диаметр окружности, точка \(N\) лежит на окружности,
\[
\angle TKN=32^\circ.
\]
Найдите:
\[
\text{1)}\ \angle TNK;
\qquad
\text{2)}\ \angle KTN.
\]
Если \(O\) --- центр окружности, найдите также
\[
\angle TON.
\]

\end{enumerate}

\subtopic{Сложный уровень}

\begin{enumerate}[resume]

\item Прямая проходит через точки
\[
A(-3;7),
\qquad
B(5;-1).
\]

\begin{enumerate}[label=\arabic*)]
\item Найдите уравнение прямой.
\item Найдите точки её пересечения с осями координат.
\item Решите неравенство
\[
y\ge2,
\]
если \(y\) задаётся найденной линейной функцией.
\item Запишите множество подходящих значений \(x\) промежутком.
\end{enumerate}

\end{enumerate}
\end{document}